\clearpage
\begingroup
\raggedbottom
\setlength{\baselineskip}{13pt}

\section{Transition Analysis and Correction}
\label{app:sentence-level-analysis}

% 完整诊断转换审计
\subsection{Full Diagnostic Transition Audit}

% 对每种监督目标，配对 11,364 个 DS 和 RF 预测。
% 后续完整诊断审计纳入全部 1,135 个有害或有益的极端转换，
% 其中 1,119 个达到多数共识。其余 16 个因三次重试后仍有一个
% 主裁判缺少有效响应，且剩余两个有效判断不一致，被记为
% no_majority 并排除。这是确定性的完成标准，而非随机子采样。
For each supervision target, we pair 11,364 DS and RF predictions. The later
full diagnostic audit includes all 1,135 harmful or beneficial extreme
transitions and reaches majority consensus on 1,119. The remaining 16 are
excluded as \textit{no\_majority}: after three retries, one primary judge still
lacks a valid response, and the two remaining valid judgments disagree. This
is a deterministic completion criterion, not random subsampling.

% 完成共识的集合包含 701 个有害转换（DS 正确而 RF 严重错误）
% 和 418 个有益转换（DS 严重错误而 RF 正确）。
% 裁判接收任务说明、评分标准、被评估文本、真实分数、
% DS/RF 预测及 RF 理由，但不知道训练条件和模型身份。
% 主裁判为 GLM-5.3-Flash 和 Doubao-Seed-2.0-Lite，
% 仲裁裁判为 MiniMax-M3。表中报告预测变化。
The completed set contains 701 harmful transitions (DS correct and RF a severe
error) and 418 beneficial transitions (DS a severe error and RF correct).
Judges receive the task instruction, rubric, evaluated text, gold score, DS and
RF predictions, and the RF rationale. The panel is GLM-5.3-Flash and
Doubao-Seed-2.0-Lite, with MiniMax-M3 as adjudicator; training condition and
model identity are hidden. Table~\ref{tab:full-interface-transitions} reports
the prediction changes.

% =========================================================
% 独立跨栏表：接口转换，表体按内容自然确定宽度
% =========================================================
\begin{table*}[!tp]
\centering
\captionsetup{font=normalsize,skip=2pt}
\small
\setlength{\tabcolsep}{3pt}
\renewcommand{\arraystretch}{1.05}

\begin{tabular}{@{}llrrrrrrr@{}}
\toprule
Training & Task
& $\Delta$QWK
& $\Delta$Acc.
& $\Delta$Adj.
& $\Delta$Sev.
& $\Delta$Inv.
& Harm.
& Bene. \\
\midrule

\multirow{4}{*}{\textbf{SO}}
& Actionability
& $-0.040$ & $-2.17$ & $-4.03$ & $+6.03$ & $+0.17$
& 130 & 46 \\

& Grounding Spec.
& $-0.018$ & $-7.00$ & $+2.97$ & $+4.00$ & $+0.03$
& 212 & 95 \\

& Helpfulness
& $-0.017$ & $-0.40$ & $-0.33$ & $+0.67$ & $+0.07$
& 0 & 0 \\

& Verifiability
& $-0.064$ & $-4.53$ & $-0.68$ & $+5.20$ & $0.00$
& 35 & 10 \\

\cmidrule(l){2-9}
\textbf{Mean} &
& $\mathbf{-0.035}$
& $\mathbf{-3.52}$
& $\mathbf{-0.52}$
& $\mathbf{+3.98}$
& $\mathbf{+0.07}$
& \textbf{377}
& \textbf{151} \\

\midrule

\multirow{4}{*}{\textbf{RS}}
& Actionability
& $-0.047$ & $-3.13$ & $\mathbf{+1.90}$ & $+2.17$
& $\mathbf{-0.93}$ & 170 & \textbf{130} \\

& Grounding Spec.
& $-0.014$ & $\mathbf{+11.10}$ & $\mathbf{-16.60}$
& $+5.50$ & $0.00$ & 128 & 128 \\

& Helpfulness
& $-0.008$ & $-0.07$ & $-0.33$ & $+0.40$ & $0.00$
& 4 & 4 \\

& Verifiability
& $\mathbf{-0.002}$ & $\mathbf{+6.05}$ & $\mathbf{-15.02}$
& $+8.97$ & $0.00$ & 36 & 7 \\

\cmidrule(l){2-9}
\textbf{Mean} &
& $\mathbf{-0.017}$
& $\mathbf{+3.49}$
& $\mathbf{-7.51}$
& $\mathbf{+4.26}$
& $\mathbf{-0.23}$
& \textbf{338}
& \textbf{269} \\

\bottomrule
\end{tabular}

% 四个有序评分任务从 DS 切换到 RF 后的逐任务变化。
% SO 和 RS 表示训练方法。准确率及各比例的变化以百分点（pp）表示。
% Harmful 表示正确预测变为严重错误；Beneficial 表示严重错误变为正确预测。
% 计数为三个随机种子上的事件总数。Mean 行加粗。
% 任务行中的粗体标记与对应 SO 单元格在符号或幅度上存在差异的 RS 结果，
% 即两种训练目标对接口切换呈现不同响应的位置。
\caption{Task-level changes after switching the inference interface from
DS to RF on the four ordinal tasks. SO and RS denote training methods.
Accuracy and ratio changes are in percentage points (pp). Harmful means
correct-to-severe-error; beneficial means severe-error-to-correct. Counts
are seed-level events across the three seeds. Mean rows are shown in bold.
In the task rows, boldface marks RS cells that differ from the corresponding
SO cells in sign or magnitude, i.e., where the two training targets respond
differently to the interface switch.}
\label{tab:full-interface-transitions}
\end{table*}

% 原始错误标签统一映射为以下展示名称：
% factual_error（事实错误）、
% evidence_misread（证据误读）、
% rubric_misapplication（评分标准误用）、
% score_mapping_error（分数映射错误）、
% unsupported_inference（无依据推断）、
% internal_contradiction（内部矛盾）、
% irrelevant_or_missing_reasoning（无关或缺失的推理）、
% other（其他）。
% “未标记错误”和“未解决的错误类型平票”属于聚合结果，
% 并不分别等同于经过验证的正确类别和实质性错误类别。
% 对应表格明确说明各自的统计分母。
Raw error labels are displayed consistently as follows:
\textit{factual\_error} (Factual error),
\textit{evidence\_misread} (Evidence misread),
\textit{rubric\_misapplication} (Rubric misapplication),
\textit{score\_mapping\_error} (Score mapping error),
\textit{unsupported\_inference} (Unsupported inference),
\textit{internal\_contradiction} (Internal contradiction),
\textit{irrelevant\_or\_missing\_reasoning} (Irrelevant or missing reasoning),
and \textit{other} (Other).
``No error flagged'' and ``Unresolved error-type tie'' are aggregation outcomes, not
validated-correct and substantive-error categories. Tables~\ref{tab:full-sentence-errors}
and~\ref{tab:full-first-error-position} state their denominators explicitly.

% =========================================================
% 独立跨栏表：句子错误标签，表体按内容自然确定宽度
% =========================================================
% SO 和 RS 分列展示，各单元格报告句子数量及其在对应训练条件下的比例。
% Delta pp 表示 RS 百分比减去 SO 百分比。
% 保留原来的标签名称、统计数值和加粗位置。
\begin{table*}[!tp]
\centering
\captionsetup{font=normalsize,skip=2pt}
\small
\setlength{\tabcolsep}{6pt}
\renewcommand{\arraystretch}{1.05}

\begin{tabular}{@{}lrrr@{}}
\toprule
Sentence label
& SO Count (\%)
& RS Count (\%)
& $\Delta$pp (RS$-$SO) \\
\midrule

No error flagged
& 687 (41.5\%)
& 657 (39.2\%)
& $-2.4$ \\

Unresolved error-type tie
& 448 (27.1\%)
& 424 (25.3\%)
& $-1.8$ \\

Rubric misapplication
& 191 (11.5\%)
& 231 (13.8\%)
& $\mathbf{+2.2}$ \\

Factual error
& 110 (6.7\%)
& 66 (3.9\%)
& $\mathbf{-2.7}$ \\

Evidence misread
& 90 (5.4\%)
& 85 (5.1\%)
& $-0.4$ \\

Score mapping error
& 72 (4.4\%)
& 132 (7.9\%)
& $\mathbf{+3.5}$ \\

Unsupported inference
& 45 (2.7\%)
& 76 (4.5\%)
& $\mathbf{+1.8}$ \\

Internal contradiction
& 10 (0.6\%)
& 5 (0.3\%)
& $-0.3$ \\

Irrelevant or missing reasoning
& 1 (0.1\%)
& 1 (0.1\%)
& $0.0$ \\

\midrule
\textbf{Total sentences}
& \textbf{1,654 (100\%)}
& \textbf{1,677 (100\%)}
& -- \\
\bottomrule
\end{tabular}

% 701 个有害事件中共 3,331 个理由句子的标签分布。
% SO 和 RS 两列分别报告句子数量，括号内为对应条件下的百分比。
% Delta pp 为 RS 百分比减去 SO 百分比，以未四舍五入的比例计算。
% 未解决的错误类型平票保留在分母中，但不纳入对应图中的错误类型分析。
\caption{Sentence labels across 3,331 rationale sentences in 701
harmful events. SO and RS columns report sentence counts with within-condition
percentages in parentheses. $\Delta$pp is the RS percentage minus the SO
percentage, computed from unrounded proportions. Unresolved ties remain
in the denominator but are excluded from
Figure~\ref{fig:rationale-error-analysis}.}
\label{tab:full-sentence-errors}
\end{table*}

% =========================================================
% 独立跨栏表：首次错误位置，表体按内容自然确定宽度
% =========================================================
\begin{table*}[!tp]
\centering
\captionsetup{font=normalsize,skip=2pt}
\small
\setlength{\tabcolsep}{8pt}
\renewcommand{\arraystretch}{1.05}

\begin{tabular}{@{}lrrrr@{}}
\toprule
& \multicolumn{2}{c}{\textbf{SO}}
& \multicolumn{2}{c}{\textbf{RS}} \\
\cmidrule(lr){2-3}\cmidrule(lr){4-5}
Position
& \shortstack{Sentence\\number}
& Quartile
& \shortstack{Sentence\\number}
& Quartile \\
\midrule

1st
& 20.7\% & 29.3\%
& \textbf{36.3\%} & \textbf{49.2\%} \\

2nd
& 28.3\% & 35.6\%
& 25.2\% & 36.6\% \\

3rd
& 36.4\% & 30.7\%
& 26.1\% & 11.7\% \\

4th+
& 14.7\% & 4.3\%
& 12.3\% & 2.4\% \\

1st--2nd
& 48.9\% & --
& \textbf{61.6\%} & -- \\

\bottomrule
\end{tabular}

\caption{First flagged-sentence position in the 701 harmful events.
Sentence number is absolute position; Quartile is relative position.}
\label{tab:full-first-error-position}
\end{table*}

\subsection{Rationale Correction and Fixed-Prefix Rescoring}

% 158 个纠错案例不是从上述 701 例全量审计中再抽样，而来自更早的独立分层审计。
% 该审计按 supervision-pair×task 最多抽 50 例，通常 harmful/beneficial 各 25，
% 小池全取并由另一方向补足，seed=20260825。
The 158-case correction set is not sampled from the 701 harmful cases above.
It comes from an earlier, separately stratified audit with seed 20260825. For
each supervision-pair--task cell, that audit selects at most 50 transitions,
typically 25 harmful and 25 beneficial; a small pool is exhausted and the
opposite direction fills the remaining quota. It selects 296 transitions
(175 harmful), completes consensus for 278, and contains 159 completed harmful
cases. Of these, 158 are judged to support the wrong score and are frozen for
the intervention. This earlier audit uses MiniMax-M3 and
Doubao-Seed-2.0-Lite as primaries and Doubao-Seed-2.1-Turbo as adjudicator,
whereas the full audit above uses the later panel and prompt version.
Only 155 of the 158 frozen record IDs overlap the later eligible set.

% 冻结集有 158 个唯一 record_id、131 个唯一 task/sample 对；允许跨 seed 或监督目标
% 重复同一测试样本。任务配额是 47/43/4/64，RS/SO 为 86/72，seed 为 56/54/48。
The frozen set has 158 unique record IDs but 131 unique task--sample pairs:
24 test samples recur across seeds or supervision targets, accounting for 51
records. The task counts are 47/43/4/64 for Actionability, Grounding
Specificity, Helpfulness, and Verifiability; RS/SO contribute 86/72 records,
and seeds 42/43/44 contribute 56/54/48. There is no separate correction-stage
quota beyond the earlier audit's stratification.

% =========================================================
% 独立跨栏表：按任务和训练目标统计恢复率
% =========================================================
\begin{table*}[!tp]
\centering
\captionsetup{font=normalsize,skip=2pt}
\small
\setlength{\tabcolsep}{8pt}
\renewcommand{\arraystretch}{1.05}

\begin{tabular}{@{}lrrr@{}}
\toprule
Task & RS & SO & All \\
\midrule

Actionability
& 21/22 & 24/25 & 45/47 \\

Grounding Spec.
& 24/25 & 18/18 & 42/43 \\

Helpfulness
& 4/4 & -- & 4/4 \\

Verifiability
& \textbf{26/35} & 25/29 & 51/64 \\

\midrule
Overall
& 75/86 & 67/72 & \textbf{142/158} \\
\bottomrule
\end{tabular}

\caption{Gold-score recovery by task and training target. Cells show
recovered/total harmful cases; Verifiability has the lowest observed rate.}
\label{tab:full-correction-task}
\end{table*}

% 保留原来的说明文字
\begingroup
\fontsize{10pt}{11.2pt}\selectfont

% 接口转换证据。
% 这些表连接接口变化、理由错误和固定前缀纠错；
% 附录中定义各项统计的分母。
\noindent\textbf{Interface-transition evidence.}
Tables~\ref{tab:full-interface-transitions}--\ref{tab:full-correction-transitions}
connect interface changes, rationale errors, and fixed-prefix correction;
Appendix~\ref{app:sentence-level-analysis} defines the denominators.
\par
\endgroup

% GLM 编辑、MiniMax 复核，最多三轮；每次调用最多三次传输/格式重试。
% 142/14/2 例分别在第 1/2/3 轮通过，共 18 次复核拒绝，最终 0 失败。
GLM-5.3-Flash edits each rationale and MiniMax-M3 independently reviews it,
for at most three edit--review rounds; each API call permits up to three
transport or format retries. Of the 158 cases, 142, 14, and 2 are approved in
rounds 1, 2, and 3, respectively, implying 18 review rejections followed by
revision and no final failures. Four records require a retried API call
(six failed call attempts in total). Rejections concern incomplete target
correction, failure to preserve valid content, inconsistent change logs,
incoherence, or wording that indirectly leaks a rubric tier.

% 编辑器和复核器不直接接收 gold/prediction 字段，但 error targets 与指导来自 gold-aware 上游审计，
% 因此流程不是端到端 gold-blind。编辑和复核先于重评分，之后不依据重评分结果继续修改。
The editor receives the task request, rubric, evaluated answer, original
rationale, and judge-derived \texttt{error\_targets} with
\texttt{correction\_guidance}. The reviewer approves only a complete,
coherent rationale that corrects every target, preserves unrelated valid
reasoning, invents no evidence, states no numeric final score or gold label,
and can be inserted for rescoring. Gold and prediction fields are not directly
sent to either model; however, the edit guidance is produced by an upstream
gold-aware audit. The workflow is therefore not end-to-end gold-blind, and
the results are evidence about rationale intervention on selected failures
rather than an estimate for the full test set. Editing and review finish
before Qwen rescoring; no rationale is revised in response to the downstream
rescoring result.

% 固定前缀重评分保留原 system/user prompt，只拼接原始或修正理由及 score 起始标签，
% 再由同一 adapter 和贪心配置续写分数。
For rescoring, the original system and user messages are rendered unchanged
with the original Qwen chat template. The only appended assistant prefix is
\texttt{<reasoning>\{rationale\}</reasoning>}\\
\texttt{\textbackslash n<score>}, using either the
original or corrected rationale; the same adapter and greedy decoding settings
then generate only the score. FixOrig reproduces NatRF in 157/158 cases.

% =========================================================
% 独立跨栏表：固定前缀重评分汇总
% =========================================================
\begin{table*}[!tp]
\centering
\captionsetup{font=normalsize,skip=2pt}
\small
\setlength{\tabcolsep}{8pt}
\renewcommand{\arraystretch}{1.05}

\begin{tabular}{@{}lrrr@{}}
\toprule
Measure & NatRF & FixOrig & FixCorr \\
\midrule
Cases & 158 & 158 & 158 \\
MAE & 2.038 & 2.044 & \textbf{0.101} \\
Incorrect & 158 & 158 & \textbf{16} \\
Returned to gold & 0 & 0 & \textbf{142 (89.9\%)} \\
\bottomrule
\end{tabular}

\caption{Aggregate fixed-prefix intervention on the frozen 158-case
set. FixOrig and FixCorr fix the original and corrected rationale,
respectively; all 474 outputs are format-valid.}
\label{tab:full-correction-aggregate}
\end{table*}

% 接续原来的分析集合说明，文字不变
Table~\ref{tab:correction-analysis-sets} distinguishes the three reported
analysis sets. The two post hoc flags are Actionability edits whose treatment
of a direct-action question conflicts with the task rubric.

% =========================================================
% 独立跨栏表：分析集合及其分母
% =========================================================
\begin{table*}[!tp]
\centering
\captionsetup{font=normalsize,skip=2pt}
\small
\setlength{\tabcolsep}{8pt}
\renewcommand{\arraystretch}{1.05}

\begin{tabular}{@{}lrr@{}}
\toprule
Analysis set & Gold recovery & Statistic \\
\midrule

Frozen intervention set
& 142/158 (89.87\%)
& MAE 0.101 \\

Excluding two post hoc rubric flags
& 140/156 (89.74\%)
& MAE 0.103 \\

FixOrig-matched subset
& 141/157 (89.81\%)
& -- \\

\bottomrule
\end{tabular}

\caption{Correction-analysis denominators. The second row excludes
\texttt{actionability\_test\_0061} and \texttt{actionability\_test\_0344};
the third uses the distinct FixOrig = NatRF restriction.}
\label{tab:correction-analysis-sets}
\end{table*}

% =========================================================
% 独立跨栏表：纠错前后的分数转换
% 两组列属于同一张表，保留原有内容结构
% =========================================================
\begin{table*}[!tp]
\centering
\captionsetup{font=normalsize,skip=2pt}
\small
\setlength{\tabcolsep}{8pt}
\renewcommand{\arraystretch}{1.05}

\begin{tabular}{@{}lrlr@{}}
\toprule
\multicolumn{2}{@{}l}{Returned to gold}
& \multicolumn{2}{l@{}}{Adjacent errors} \\
\midrule
$5\to3$ & 63 & $1\to2$ (gold 3) & 11 \\
$1\to3$ & 30 & $5\to2$ (gold 3) & 3 \\
$3\to5$ & 22 & $3\to2$ (gold 1) & 1 \\
$3\to1$ & 13 & $5\to4$ (gold 3) & 1 \\
$4\to2$ & 4 & & \\
$5\to2$ & 3 & & \\
$2\to4$ & 3 & & \\
$4\to1$ & 3 & & \\
$2\to5$ & 1 & & \\
\midrule
Subtotal & 142 & Subtotal & 16 \\
\bottomrule
\end{tabular}

\caption{FixOrig-to-FixCorr score transitions. All 158 predictions
move closer to gold: 142 return exactly and 16 become adjacent errors, reducing
total absolute error from 323 to 16.}
\label{tab:full-correction-transitions}
\end{table*}

\endgroup